\documentclass[11pt]{article}

\usepackage[margin=1in]{geometry}
\usepackage{amsmath,amssymb,amsthm,mathtools}
\usepackage{booktabs}
\usepackage{enumitem}
\usepackage[colorlinks=true,linkcolor=blue,citecolor=blue,urlcolor=blue]{hyperref}

\newtheorem{theorem}{Theorem}[section]
\newtheorem{proposition}[theorem]{Proposition}
\newtheorem{lemma}[theorem]{Lemma}
\newtheorem{corollary}[theorem]{Corollary}
\theoremstyle{definition}
\newtheorem{remark}[theorem]{Remark}
\newtheorem{observation}[theorem]{Observation}

\newcommand{\R}{\mathbb R}
\newcommand{\Q}{\mathbb Q}
\newcommand{\E}{\mathbb E}
\newcommand{\Var}{\operatorname{Var}}
\newcommand{\Cov}{\operatorname{Cov}}
\newcommand{\Tr}{\operatorname{Tr}}
\newcommand{\tw}{T_W}
\newcommand{\tu}{T_U}
\newcommand{\kw}{K_W}
\newcommand{\thet}{\theta_{1,2,4}}

\title{De-flagging the smoothed Goodman certificate\\
\large What the $C_3\cup_{K_2}C_5$ proof says without flag algebra,
and what a human proof still has to beat}
\author{Working note}
\date{August 20, 2026}

\begin{document}
\maketitle

\begin{abstract}
\texttt{smoothed\_goodman\_theorem.tex} proves $\Delta_2(W)=t(\thet,W)-(2p-1)t(C_5,W)\ge0$
by an exact rational multiplier flag-SOS certificate at ground size $8$.  This note
answers the question ``can the flag algebra be removed?''.  The answer has two
halves.  \emph{Yes as language}: the certificate is, verbatim, the statement that
$\Delta_2$ is a nonnegative combination of (i) integrals of squares of rooted
subgraph densities, (ii) $(2p-1)$ times such integrals, and (iii) induced
$8$-vertex densities; we state this in pure graphon language and prove the three
supporting lemmas from scratch, so that no flag-algebra formalism is needed to
read or trust the logic (\S\ref{sec:deflag}).  \emph{No as computation}: we
measured the committed certificate --- $170+14$ PSD blocks of total dimension
$13{,}242$ (largest $272\times272$), $1{,}431{,}052$ integer entries over the
common denominator $4^{48}=2^{96}$, and $12{,}346$ per-graph inequalities
(\S\ref{sec:size}).  Translation is a change of vocabulary, not of length.

We then record what \emph{is} human-scale.  We give a flag-free reduction chain
with two exact identities (\S\ref{sec:structure}), verified in rational
arithmetic against the paper's own gates; we kill four natural flag-free routes
with explicit witnesses, one of them exactly rational (\S\ref{sec:dead}); and we
measure how much room a lossy human argument actually has
(\S\ref{sec:landscape}).  The headline of that measurement: routed through
Goodman's inequality, the residual ``correlation deficit'' never exceeded $2\%$
of the available budget over ${\sim}10^6$ step graphons including adversarial
hill-climbing, and it is \emph{negative} near the balanced multipartite
frontier.  This corrects the expectation recorded in
\texttt{smoothed\_goodman\_spectral\_partial.tex}~\S7 that the frontier is where
the domination is tight.  The genuine obstruction is instead located precisely:
the Goodman part $G$ vanishes identically on the whole family of complete
multipartite graphons, balanced \emph{or not}, where the theorem degenerates to
a vertex-bias covariance inequality that is false for general graphons
(\S\ref{sec:landscape}, Observation~\ref{obs:multipartite}).
\end{abstract}

\tableofcontents

\section{Notation and the two statements}\label{sec:notation}

$W:[0,1]^2\to[0,1]$ is a graphon, $U:=1-W$, $p:=t(K_2,W)=\iint W$, $q:=1-p$.
$\tw$ is the integral operator with kernel $W$, and $\tw^r(x,y)$ denotes the
kernel of its $r$-th power, so $t(C_m,W)=\Tr(\tw^m)=\sum_k\lambda_k^m$ for the
eigenvalues $\lambda_k$ of $\tw$.  Write
\[
  d(x):=\int W(x,y)\,dy,\qquad a:=d-p,\qquad c_5(x):=\tw^5(x,x).
\]
$\thet=C_3\cup_{K_2}C_5$ is the theta graph with branch paths of lengths
$1,2,4$, so that
$t(\thet,W)=\iint W(x,y)\,\tw^2(x,y)\,\tw^4(x,y)\,dx\,dy$.  The target is
\[
  \Delta_2(W)\;=\;t(\thet,W)-(2p-1)\,t(C_5,W)
  \;=\;\iint W(x,y)\big(\tw^2(x,y)-(2p-1)\big)\tw^4(x,y)\,dx\,dy\;\ge\;0 .
\]
We write $\kw$ for the \emph{Goodman defect kernel}
$\kw(x,y)=W(x,y)\big(\tw^2(x,y)-(2p-1)\big)$, so $\Delta_2=\langle
\kw,\tw^4\rangle=\Tr(\tw^2\kw\tw^2)$.  Two scalars recur:
\[
  \Gamma\;:=\;t(K_3,W)-p(2p-1)\ \ge0 \quad(\text{Goodman slack}),
  \qquad
  G\;:=\;\iint W(x,y)\,\tu^2(x,y)\,\tw^4(x,y)\,dx\,dy\ \ge0 .
\]
Throughout, ``step graphon'' means a symmetric block matrix $B\in[0,1]^{n\times
n}$ with block masses $w\in\Delta_{n-1}$; all numerical statements below were
computed with an independent implementation, validated against the two exact
gates published in \texttt{smoothed\_goodman\_theorem.tex}: it reproduces
$\Delta_2(W\equiv p)=p^5(1-p)^2$ and $\Delta_2=\tfrac{177}{25000}$ at the
complete tripartite profile $(\tfrac12,\tfrac3{10},\tfrac15)$, both in exact
rational arithmetic.

\section{The certificate without flag algebra}\label{sec:deflag}

Flag algebra is a bookkeeping formalism for three elementary facts about
graphons.  We state them, prove them, and then state exactly what the committed
certificate asserts.  Nothing below mentions flags, types, or the flag product.

\subsection{Rooted densities and squares}

Fix $m\ge0$ and a graph $\sigma$ on the labelled vertex set $[m]$ (the
\emph{root pattern}).  For $\mathbf x=(x_1,\dots,x_m)\in[0,1]^m$ put
\[
  W^\sigma(\mathbf x)\;:=\;\prod_{ij\in E(\sigma)}W(x_i,x_j)
  \prod_{ij\notin E(\sigma)}\big(1-W(x_i,x_j)\big),
\]
the density of the pattern $\sigma$ at the roots.  For a graph $F$ containing
$\sigma$ as a labelled induced subgraph, with $v(F)=\ell$, let
\[
  t_\sigma(F;\mathbf x)\;:=\;\int_{[0,1]^{\ell-m}}
  \prod_{ij\in E(F)}W(y_i,y_j)\prod_{ij\notin E(F)}\big(1-W(y_i,y_j)\big)
  \,dy_{m+1}\cdots dy_\ell ,
\]
where $y_i=x_i$ for $i\le m$.  This is the \emph{rooted induced density} of $F$
at $\mathbf x$: the probability that $\ell-m$ fresh random points, together with
the roots, induce $F$.

\begin{lemma}[Squares are nonnegative; this is all ``Gram positivity'' means]
\label{lem:square}
Let $\sigma$ be a root pattern on $[m]$, let $F_1,\dots,F_n$ be graphs on
$\ell$ vertices containing $\sigma$ as a labelled induced subgraph, and let
$u\in\R^n$.  Then
\[
  \int_{[0,1]^m}W^\sigma(\mathbf x)
  \Big(\sum_{i=1}^{n}u_i\,t_\sigma(F_i;\mathbf x)\Big)^{\!2}d\mathbf x\;\ge\;0 .
\]
Consequently, for any positive semidefinite $Q\in\R^{n\times n}$,
$\int W^\sigma(\mathbf x)\,f(\mathbf x)^{T}Qf(\mathbf x)\,d\mathbf x\ge0$, where
$f(\mathbf x)=\big(t_\sigma(F_i;\mathbf x)\big)_{i\le n}$.
\end{lemma}

\begin{proof}
The integrand is a nonnegative weight times a square.  For the consequence,
write $Q=\sum_j c_j u_ju_j^{T}$ with $c_j\ge0$ and sum.
\end{proof}

The only genuinely combinatorial input is that the quantity
$\int W^\sigma f f^{T}$ --- call it $\mathrm{Gram}^\sigma(W)$ --- can be
expanded in \emph{unrooted} induced densities of graphs on $2\ell-m$ vertices.
That expansion is a finite counting identity (choose which of the $2\ell-m$
points play the roles of the roots and of the two copies of the fresh points),
and it is what produces the integer tables $M^\sigma(H)$ of the certificate:
\begin{equation}\label{eq:gram}
  \sum_{H}p_{N}(H,W)\,M^{\sigma}(H)\;=\;\mathrm{Gram}^{\sigma}(W)\;\succeq\;0,
  \qquad N=2\ell-m,
\end{equation}
where $p_N(H,W)$ is the probability that $N$ i.i.d.\ uniform points induce $H$.
Identity \eqref{eq:gram} is a statement about integers only; it is the part of
the pipeline that is machine-generated and double-derived.

\subsection{The multiplier is an independence statement}

\begin{lemma}[Factorisation]\label{lem:mult}
Let $\sigma'$ be a root pattern with $2\ell'-m'=6$, let $M^{\sigma'}_6$ be its
ground-$6$ table, and for a graph $H$ on $8$ vertices set
\[
  \mathrm{Mult}^{\sigma'}(H)\;:=\;\frac1{\binom82}\sum_{\{a,b\}\subseteq V(H)}
  \big(2A_H(a,b)-1\big)\,M^{\sigma'}_6\big(H[V(H)\setminus\{a,b\}]\big).
\]
Then for every graphon $W$,
$\sum_H p_8(H,W)\,\mathrm{Mult}^{\sigma'}(H)=(2p-1)\,\mathrm{Gram}^{\sigma'}(W)$.
\end{lemma}

\begin{proof}
Sample $x_1,\dots,x_8$ i.i.d.\ uniform and let $G$ be the induced random graph.
Averaging the definition against $p_8(\cdot,W)$ and using exchangeability gives
$\E\big[(2\mathbf 1[x_1x_2\in G]-1)\,M^{\sigma'}_6(G[x_3,\dots,x_8])\big]$.  The
edge indicator of $\{x_1,x_2\}$ depends only on $(x_1,x_2)$ and on the edge
variable of that pair; $G[x_3,\dots,x_8]$ depends only on $(x_3,\dots,x_8)$ and
on the edge variables of pairs inside $\{3,\dots,8\}$.  These are disjoint sets
of independent variables, so the expectation factorises, and
$\E[2W(x_1,x_2)-1]=2p-1$.
\end{proof}

Lemma~\ref{lem:mult} is exactly where the earlier \emph{invalid} attempts
failed: multiplying a same-ground table by the per-graph scalar $(2p_H-1)$ does
not factorise, because then the edge and the $6$-set share coordinates.  The
disjointness is the whole content.

\subsection{What the certificate asserts}

\begin{theorem}[De-flagged form of the ground-$8$ multiplier certificate]
\label{thm:deflag}
There exist finitely many root patterns $\sigma_1,\dots,\sigma_{170}$ and
$\sigma'_1,\dots,\sigma'_{14}$, positive semidefinite rational matrices
$Q_b,R_{b'}$, and rationals $s_H\ge0$ $(H\in\mathcal H_8)$, such that for every
graphon $W$,
\begin{equation}\label{eq:deflagged}
  \Delta_2(W)\;=\;
  \sum_{b=1}^{170}\int W^{\sigma_b}f_{b}^{T}Q_bf_{b}
  \;+\;(2p-1)\sum_{b'=1}^{14}\int W^{\sigma'_{b'}}f_{b'}^{T}R_{b'}f_{b'}
  \;+\;\sum_{H\in\mathcal H_8}s_H\,p_8(H,W).
\end{equation}
In words: \emph{$\Delta_2$ is a nonnegative combination of integrals of squares
of rooted densities, of $(2p-1)$ times such integrals, and of induced
$8$-vertex densities.}  In particular $\Delta_2(W)\ge0$ whenever $p\ge\tfrac12$;
and for $p\le\tfrac12$, $\Delta_2=t(\thet,W)+(1-2p)t(C_5,W)\ge0$ termwise.
\end{theorem}

\begin{proof}
Given the per-graph inequalities $s_H:=c[H]-\sum_b\langle
Q_b,M^{\sigma_b}_8(H)\rangle-\sum_{b'}\langle
R_{b'},\mathrm{Mult}^{\sigma'_{b'}}(H)\rangle\ge0$ that the verifier checks in
exact integer arithmetic, multiply by $p_8(H,W)\ge0$, sum over $H$, and use
$\Delta_2=\sum_Hp_8(H,W)c[H]$ together with \eqref{eq:gram} and
Lemma~\ref{lem:mult}.  Nonnegativity of the first two groups is
Lemma~\ref{lem:square}, using $2p-1\ge0$ for the second.
\end{proof}

\begin{remark}
Equation~\eqref{eq:deflagged} is an \emph{identity}, not merely an inequality:
the slacks $s_H$ are defined so that it holds exactly.  The certificate already
stores each $Q_b$ in factored form $Q_b=4^{-s}\,U_bF_b(I+\Theta_b)F_b^{T}U_b^{T}$,
which is literally a sum of squares once $I+\Theta_b$ is diagonalised; and
$I+\Theta_b\succ0$ is certified by the exact rational check
$\|\Theta_b\|_F<1$.  So no eigenvalue computation, and no flag-algebra
vocabulary, enters the proof at any point.
\end{remark}

\section{Why the translation stops there: the measured size}\label{sec:size}

Theorem~\ref{thm:deflag} is the honest translation.  It removes the formalism
completely, and it reduces the trust chain to ``squares are nonnegative'' plus
``disjoint coordinates are independent'' plus one integer computation.  What it
does not do is make the proof short, and it cannot: we measured the committed
certificate \texttt{flag\_sdp/n8\_mult/rounding/data/certificate.json.gz}
directly.

\begin{center}
\begin{tabular}{lr}
\toprule
ground-$8$ PSD blocks $Q_b$ & $170$\\
\quad block sizes & $7$ to $272$; $133$ of them are $64\times64$\\
ground-$6$ multiplier blocks $R_{b'}$ & $14$ (sizes $17^2,16^6,14^5,1$)\\
total PSD dimension & $13{,}041+201=13{,}242$\\
integer entries in the stored factors & $1{,}431{,}052$\\
common denominator & $4^{48}=2^{96}$\\
rational $\Theta$ corrections & $53$ entries across $8$ blocks\\
rank-one additions & $0$\\
per-graph inequalities $s_H\ge0$ & $12{,}346$\\
zero slacks & $48$ ($22$ complete multipartite, $26$ bipartite-family)\\
\bottomrule
\end{tabular}
\end{center}

The full ground-$8$ size multiset is $133$ blocks of size $64$, then
$13{\times}62$, $3{\times}61$, $3{\times}55$, $2{\times}272$, $2{\times}269$,
$2{\times}260$, $2{\times}60$, and one each of sizes
$270,267,266,265,252,109,104,63,50,7$.

Writing \eqref{eq:deflagged} out explicitly therefore means exhibiting on the
order of $10^6$ rational numbers with $96$-bit denominators, plus $12{,}346$
verified inequalities.  There is no mathematical obstruction to doing so --- it
is what the file contains --- but the result is not a proof a human reads; it is
the same computer proof in a different alphabet.  Two consequences worth
stating plainly:

\begin{enumerate}[leftmargin=1.6em]
\item \emph{The flag algebra is not the reason the proof is hard to verify.}
The reason is the $12{,}346$ inequalities and the $13{,}242$-dimensional PSD
data.  De-flagging changes what a referee must \emph{understand} (three
elementary lemmas instead of Razborov's calculus) but not what a referee must
\emph{check} (a $150$-line integer program, $5$ seconds).
\item \emph{No block is idle.} All $170$ ground-$8$ and all $14$ ground-$6$
blocks have nonzero factors, so there is no cheap sparsification of the existing
certificate into a human-scale one.  A short proof, if it exists, must come from
a different argument, not from compressing this one.
\end{enumerate}

The rest of this note is about that different argument: what flag-free structure
is available, which natural routes are provably dead, and --- the useful part
--- how much room a lossy human estimate actually has.

\section{Flag-free structure that is human-scale}\label{sec:structure}

Everything in this section is elementary and was verified in exact rational
arithmetic on step graphons.

\subsection{Two parallel decompositions}

The pointwise identity $\tw^2(x,y)-(2p-1)=\tu^2(x,y)+a(x)+a(y)$ (expand
$W=1-U$) contracts against two different weights and gives two statements that
should be read side by side.

\begin{proposition}[Goodman and its smoothed form]\label{prop:parallel}
For every graphon $W$,
\begin{align}
  \Gamma\;=\;t(K_3,W)-p(2p-1)\;&=\;\iint W\,\tu^2\;+\;2\Var(d),
  \label{eq:goodman-split}\\
  \Delta_2(W)\;&=\;\underbrace{\iint W\,\tu^2\,\tw^4}_{=\;G\;\ge\;0}
  \;+\;2\Cov\big(d,c_5\big).
  \label{eq:delta-split}
\end{align}
\end{proposition}

\begin{proof}
Contract the pointwise identity against $W(x,y)\,dx\,dy$ and against
$W(x,y)\tw^4(x,y)\,dx\,dy$ respectively.  In both cases the two degree terms are
equal by symmetry.  For \eqref{eq:goodman-split},
$\iint W(x,y)a(x)=\int a\,d=\Var(d)$ since $\int a=0$.  For
\eqref{eq:delta-split}, $\int_yW(x,y)\tw^4(x,y)\,dy=\tw^5(x,x)=c_5(x)$, and
$\int a\,c_5=\Cov(d,c_5)$ since $\int a=0$.  Nonnegativity of $G$ holds because
$W\ge0$, $\tu^2\ge0$ and $\tw^4\ge0$ pointwise.
\end{proof}

Thus Goodman's inequality is the case ``weight $\equiv1$'' and the theorem is
the case ``weight $=\tw^4$'' of one statement.  Goodman is easy only because in
\eqref{eq:goodman-split} the degree term is a \emph{variance}, hence
nonnegative; in \eqref{eq:delta-split} it is a \emph{covariance}, and
$\Cov(d,c_5)<0$ does occur.  That single sign change is the entire difficulty.

\subsection{An exact spectral refinement}

Expanding the covariance in the eigenbasis $\{\varphi_k\}$ of $\tw$ isolates the
negative part of the spectrum as the sole source of trouble.

\begin{proposition}[Dirichlet form of $\Delta_2$]\label{prop:dirichlet}
Let $E_k:=\tfrac12\iint W(x,y)\big(\varphi_k(x)-\varphi_k(y)\big)^2dx\,dy\ \ge0$
be the Dirichlet energy of the $k$-th eigenfunction.  Then
\begin{equation}\label{eq:dirichlet}
  \Delta_2(W)\;=\;G\;+\;2\big(t(C_6,W)-p\,t(C_5,W)\big)\;+\;2\sum_k\lambda_k^5E_k .
\end{equation}
\end{proposition}

\begin{proof}
$\Cov(d,c_5)=\sum_k\lambda_k^5\big(\int d\varphi_k^2-p\big)$ because
$c_5=\sum_k\lambda_k^5\varphi_k^2$ and $\int c_5=t(C_5)$.  The Dirichlet
identity $\int d\,\varphi_k^2-\langle\varphi_k,\tw\varphi_k\rangle=E_k$ gives
$\int d\varphi_k^2=\lambda_k+E_k$, so
$\Cov(d,c_5)=\sum_k\lambda_k^5(\lambda_k-p+E_k)
=t(C_6)-p\,t(C_5)+\sum_k\lambda_k^5E_k$.  Substitute into
\eqref{eq:delta-split}.
\end{proof}

Three immediate readings of \eqref{eq:dirichlet}:

\begin{enumerate}[leftmargin=1.6em]
\item If $W$ is regular then $a\equiv0$, and \eqref{eq:delta-split} already
gives $\Delta_2=G\ge0$; no spectral information is needed.
\item If $\tw\succeq0$ then every $\lambda_k^5E_k\ge0$, so
$\Delta_2\ge2\big(t(C_6)-p\,t(C_5)\big)$; the PSD case would follow from
$t(C_6)\ge p\,t(C_5)$.  That inequality is false --- see
\S\ref{sec:dead}(d) --- but it is the cleanest sufficient condition we found for
the PSD case, and it is considerably simpler than the crux
$(\mathrm{JEN})$/$(\mathbf A)$ of \texttt{smoothed\_goodman\_spectral\_partial.tex}.
\item At the balanced $m$-partite frontier $W_m$ the three terms are
$0$, $+\tfrac{m-1}{m^5}$ and $-\tfrac{m-1}{m^5}$: the equality case is an exact
cancellation between the cycle comparison and the negative-eigenvalue Dirichlet
term, not a termwise vanishing.
\end{enumerate}

Identity \eqref{eq:dirichlet} was checked against a direct evaluation of
$\Delta_2$ on $6\times10^4$ random step graphons ($n\le7$ blocks); the largest
absolute discrepancy was $2.9\times10^{-15}$.

\subsection{The Goodman chain}

Dividing \eqref{eq:delta-split} differently gives the form in which the theorem
has the most visible slack.  Let $\nu:=p^{-1}W(x,y)\,dx\,dy$ be the edge
probability measure and put
\[
  \Lambda\;:=\;\frac{t(K_3,W)\,t(C_5,W)}{p}-t(\thet,W)
  \;=\;-\,p\,\Cov_\nu\!\big(\tw^2,\tw^4\big).
\]

\begin{proposition}[Chain identity]\label{prop:chain}
For every graphon $W$ with $p>0$,
\begin{equation}\label{eq:chain}
  \Delta_2(W)\;=\;\frac{t(C_5,W)}{p}\,\Gamma\;-\;\Lambda .
\end{equation}
\end{proposition}

\begin{proof}
$\frac{t(C_5)}{p}\Gamma-\Lambda
=\frac{t(C_5)t(K_3)}{p}-(2p-1)t(C_5)-\frac{t(K_3)t(C_5)}{p}+t(\thet)
=t(\thet)-(2p-1)t(C_5)$.
\end{proof}

So the theorem says exactly: \emph{the correlation deficit $\Lambda$ between
codegree and $4$-walk density, over a random edge, never exceeds the Goodman
slack budget $t(C_5)\Gamma/p$.}  Both sides vanish at every equality case.  If
$\Lambda\le0$ --- that is, if $\tw^2$ and $\tw^4$ were positively correlated
over a random edge --- the theorem would follow from Goodman alone.  They are
not always positively correlated (\S\ref{sec:dead}(a)), but \S\ref{sec:landscape}
shows the failure is small by a factor of about $50$.

\section{Four natural flag-free routes, and why each fails}\label{sec:dead}

Each of the following would give a short proof.  Each is false.  We record
explicit witnesses so that the routes need not be re-explored.

\paragraph{(a) Positive correlation of $\tw^2$ and $\tw^4$ over a random edge.}
The statement $\Lambda\le0$, i.e.
\[
  (\mathrm{CORR})\qquad p\cdot t(\thet,W)\;\ge\;t(K_3,W)\,t(C_5,W),
\]
would, with \eqref{eq:chain} and Goodman, prove $\Delta_2\ge0$ for \emph{all}
$p$ in two lines.  It is false, already on two blocks and already on the top
branch.  Take the step graphon
\[
  B=\begin{pmatrix}3/10&1\\1&0\end{pmatrix},\qquad w=\big(\tfrac45,\tfrac15\big),
  \qquad p=\tfrac{64}{125}\;>\;\tfrac12 .
\]
Then, exactly,
$t(K_3)=\tfrac{2016}{15625}$, $t(C_5)=\tfrac{415776}{9765625}$,
$t(\thet)=\tfrac{13003584}{1220703125}$, and
\[
  p\,t(\thet)-t(K_3)\,t(C_5)\;=\;-\,\frac{1195008}{30517578125}
  \;=\;-3.9158\ldots\times10^{-5}\;<\;0
\]
(a relative deficit of $0.36\%$), while
$\Delta_2=\tfrac{11756256}{1220703125}>0$, as it must be.  This is the cleanest
obstruction in the note: the theorem is \emph{not} a corollary of Goodman plus a
correlation inequality.

\paragraph{(b) Copositivity of the Goodman defect kernel.}
Since $\Delta_2=\langle\kw,\tw^4\rangle$ and $\tw^4$ is both positive
semidefinite and pointwise nonnegative, it would suffice that $\kw$ be
copositive, i.e.\ $\langle v,\kw v\rangle\ge0$ for every $v\ge0$.  It is not,
even for $p$ close to $1$.  Let $\varepsilon\in(0,\tfrac13)$ and take three
blocks $S_1,S_2,R$ of masses $\varepsilon,\varepsilon,1-2\varepsilon$ with
\[
  B=\begin{pmatrix}0&1&0\\1&0&0\\0&0&1\end{pmatrix},
\]
that is: $S_1$ and $S_2$ are completely joined to each other and to nothing
else, and $R$ is a clique.  Then $x\in S_1,y\in S_2$ have no common neighbour,
so $\kw(x,y)=-(2p-1)<0$, and with $v=\mathbf 1_{S_1\cup S_2}$,
\[
  \langle v,\kw v\rangle=-2\varepsilon^2(2p-1)<0 .
\]
Numerically: $\varepsilon=0.1$ gives $p=0.66$ and
$\langle v,\kw v\rangle=-6.4\times10^{-3}$, while $\Delta_2=0.157>0$.  So the
positivity of $\Delta_2$ genuinely uses that the test kernel is $\tw^4$ and not
an arbitrary nonnegative one.

\paragraph{(c) Rooting the pentagon.}
Writing $\tw^4(x,y)=\int\tw^2(x,z)\tw^2(z,y)\,dz$ gives
$\Delta_2=\int_z\langle v_z,\kw v_z\rangle\,dz$ with $v_z:=\tw^2(z,\cdot)\ge0$,
which suggests proving the integrand nonnegative for each root $z$.  That fails
too, for the same reason as (b) but with a witness inside the range of $\tw^2$:
\[
  B=\begin{pmatrix}1&0&0\\0&1&1\\0&1&0\end{pmatrix},\qquad
  w=(0.0721,\,0.4189,\,0.5091),\qquad p=0.6071,
\]
a graphon with a small clique component $S$ of mass $0.0721$.  Rooting at
$z\in S$ gives $v_z=0.0721\cdot\mathbf 1_S$, supported where the internal
codegree $0.0721$ falls below $2p-1=0.2142$, so
$\langle v_z,\kw v_z\rangle=-3.83\times10^{-6}<0$.  The mechanism is generic:
any small dense component has internal codegree below the global $2p-1$.

\paragraph{(d) The cycle comparison $t(C_6)\ge p\,t(C_5)$.}
By Proposition~\ref{prop:dirichlet} this would close the positive semidefinite
case outright.  It is false, both for general graphons (worst normalised deficit
$-0.034$ at $p\approx0.39$; $-0.025$ at $p\approx0.54$ on the top branch) and,
more importantly, for PSD graphons ($-0.025$ at $p\approx0.26$; $-0.021$ at
$p\approx0.51$ on the top branch), over $2\times10^5$ sampled step graphons of
each kind with $n\le8$ blocks.  So the PSD case does not collapse to a pure
spectral comparison, consistent with the earlier finding that it resists
majorisation arguments.

\paragraph{(e) Cauchy--Schwarz on the covariance.}  Both decompositions
\eqref{eq:delta-split} and \eqref{eq:chain} invite bounding the covariance by
the product of standard deviations.  Both are far too lossy.  On the chain,
$\sqrt{\Var_\nu(\tw^2)\Var_\nu(\tw^4)}\big/\big(\E_\nu[\tw^4]\cdot
\E_\nu[\tw^2-(2p-1)]\big)$ must be $\le1$; its measured maximum is $14.6$.  On
\eqref{eq:delta-split}, $2\sqrt{\Var(d)\Var(c_5)}/G$ must be $\le1$; its
measured maximum is $9.2\times10^{4}$.  The reason for the second blow-up is
structural and is the subject of Observation~\ref{obs:multipartite} below.

\section{How much room a human proof has}\label{sec:landscape}

The two decompositions turn the theorem into a ratio being at most $1$:
\[
  R\;:=\;\frac{\Lambda}{\,t(C_5)\Gamma/p\,}\ \le\ 1
  \quad\Longleftrightarrow\quad \Delta_2\ge0
  \quad\Longleftrightarrow\quad
  \rho\;:=\;\frac{-2\Cov(d,c_5)}{G}\ \le\ 1 .
\]
Neither ratio is a new statement --- each is the theorem rewritten --- but their
\emph{numerical size} tells a would-be prover how lossy an estimate may be.  We
measured both over $p\ge\tfrac12$ step graphons: ${\sim}10^6$ random samples
($n\le7$ blocks, three samplers: uniform, Bernoulli, clipped-affine), plus
hill-climbing with $60$ restarts per block count $n=2,\dots,6$ directly
maximising $R$.

\begin{center}
\begin{tabular}{lcc}
\toprule
& $\max R$ & $\max\rho$\\
\midrule
random sampling, $p\ge\tfrac12$        & $0.0138$ & $0.0671$\\
adversarial hill-climbing, $n\le6$     & $0.0198$ & ---\\
perturbations of balanced $W_m$, $m=3,4,5$, $\varepsilon\to0$
                                       & $<0$ (to $-0.06$) & $<0$ (to $-0.63$)\\
\bottomrule
\end{tabular}
\end{center}

Two conclusions.

\paragraph{The frontier is safe, not tight.}
\texttt{smoothed\_goodman\_spectral\_partial.tex}~\S7 records the expectation
that $G\ge2D$ is ``frontier-tight'', with $2D/G\to1$ as $W\to W_m$.  The
measurements do not support this.  Perturbing $W_m$ ($m=3,4,5$) in both the
edge-weight and the part-mass directions, with $\varepsilon$ from $10^{-1}$ down
to $10^{-5}$, the ratio $\rho$ is \emph{negative} throughout and tends to $0$
from below (at $m=4$: mean $-0.74,-0.099,-0.011,-0.0009,-0.0001$ for
$\varepsilon=10^{-1},\dots,10^{-5}$).  The degree correction \emph{helps} near
the frontier rather than competing with $G$; this is consistent with the
first-order mechanism of the quantitative neighbourhood theorem there, and it
means the frontier is not where a proof will be lost.

\paragraph{The real obstruction is the complete multipartite family.}

\begin{observation}[$G$ vanishes on all complete multipartite graphons]
\label{obs:multipartite}
Let $W$ be complete multipartite with arbitrary part masses.  Then $U$ is
supported on the union of the diagonal blocks, so for $x,y$ in different parts
$\tu^2(x,y)=\int U(x,z)U(z,y)\,dz=0$: no $z$ lies in both parts.  Hence
$\tu^2\equiv0$ on $\operatorname{supp}W$ and $G=0$ identically, so by
\eqref{eq:delta-split}
\[
  \Delta_2(W)\;=\;2\Cov\big(d,c_5\big)\qquad\text{on this whole family.}
\]
\end{observation}

We verified this exactly at the unbalanced tripartite profile
$(\tfrac12,\tfrac3{10},\tfrac15)$: there $G=0$ and
$2\Cov(d,c_5)=\Delta_2=\tfrac{177}{25000}$, reproducing the value published in
\texttt{smoothed\_goodman\_theorem.tex}, Prop.~5.1.  Also
$\Gamma=\tfrac{39}{1250}$ and the split \eqref{eq:goodman-split} holds exactly.

This is the sharp statement of where the difficulty lives.  On a
$(k-1)$-parameter family of extremal-looking graphons the Goodman part $G$ is
identically zero, so the decomposition \eqref{eq:delta-split} offers \emph{no}
budget at all, and the theorem there is exactly the vertex-bias inequality
\[
  \Cov\big(d,c_5\big)\;\ge\;0,
\]
which is false for general graphons (this is the ``naive vertex-bias
inequality'' already known to fail).  So a proof via \eqref{eq:delta-split} must
prove $\Cov(d,c_5)\ge0$ on the multipartite family by a special argument and
then perturb --- which is why the generic $15$-fold slack in $\rho$ is
misleading, and why Cauchy--Schwarz blows up by $10^5$ there.

By contrast the chain \eqref{eq:chain} does \emph{not} degenerate on that
family: $\Gamma=\tfrac{39}{1250}>0$ and $t(C_5)>0$ at the unbalanced profile, so
the budget $t(C_5)\Gamma/p$ is strictly positive wherever $\Delta_2>0$.  This
makes \eqref{eq:chain} the better target.

\section{Verdict and recommendation}\label{sec:verdict}

\begin{enumerate}[leftmargin=1.6em]
\item The flag algebra \emph{can} be removed from the existing proof, and
Theorem~\ref{thm:deflag} does so: the certificate asserts nothing more exotic
than ``a nonnegative combination of squares of rooted densities, of $(2p-1)$
times such squares, and of induced densities''.  Lemmas~\ref{lem:square} and
\ref{lem:mult} are the only conceptual content, and both are three-line proofs.
A referee who distrusts flag algebra can read Theorem~\ref{thm:deflag} instead
and lose nothing.
\item The computation \emph{cannot} be removed: $13{,}242$ PSD dimensions,
$1.4\times10^6$ rationals over $2^{96}$, $12{,}346$ inequalities, no idle
blocks.  ``Verified easily'' will not be achieved by translation; it would be
achieved by formalising the $150$-line integer verifier, which is a realistic
Lean target and, given the exact-rational structure already in place, probably
the cheapest route to a machine-checked proof of the theorem as it stands.
\item For a genuine human proof, the target to attack is \eqref{eq:chain}:
prove $\Lambda\le t(C_5)\Gamma/p$, i.e.\ that the edge-correlation deficit
between codegree and $4$-walk density is dominated by the Goodman slack.  The
measured true ratio never exceeded $0.02$, so an argument may be lossy by a
factor of $50$ and still succeed --- but it must be sign-aware: plain
Cauchy--Schwarz overshoots by $14.6$.  The witness in \S\ref{sec:dead}(a) shows
the deficit is genuinely positive sometimes, so the estimate cannot be a pure
correlation inequality.
\item Do not attack via \eqref{eq:delta-split} alone:
Observation~\ref{obs:multipartite} shows its budget vanishes identically on the
complete multipartite family, where the residual statement
$\Cov(d,c_5)\ge0$ is false in general and must be proved by other means.
\end{enumerate}

\begin{remark}[Scope]
All numerical statements in \S\ref{sec:dead} and \S\ref{sec:landscape} are
searches over step graphons with at most $7$--$8$ blocks; they are evidence
about sizes of ratios, not proofs.  The exact statements --- the witnesses in
\S\ref{sec:dead}(a),(b),(c), Observation~\ref{obs:multipartite}, and the
identities of \S\ref{sec:structure} --- are proved or verified in rational
arithmetic and stand on their own.
\end{remark}

\end{document}
